\documentclass[a4paper]{arrowhead}

\usepackage[yyyymmdd]{datetime}
\usepackage{etoolbox}
\usepackage[T1]{fontenc}
\usepackage[utf8]{inputenc}
\usepackage{multirow}
\usepackage{longtable}
\usepackage{amssymb}
\usepackage{float}
\usepackage{tikz}
\usetikzlibrary{positioning,arrows.meta,fit,calc,shapes.misc,shapes.geometric}

\renewcommand{\dateseparator}{-}
\setlength{\emergencystretch}{3em}
\lstset{basicstyle=\ttfamily\footnotesize,breaklines=true,columns=fullflexible,frame=single,rulecolor=\color{gray!50},backgroundcolor=\color{gray!6}}
\tikzset{
  ahsys/.style={draw=ArrowheadBlue,very thick,fill=ArrowheadBlue!6,rounded corners=2pt,minimum width=3.2cm,minimum height=1.1cm,align=center,font=\small\bfseries},
  ahext/.style={draw=ArrowheadGrey,thick,dashed,fill=gray!5,rounded corners=2pt,minimum width=3.2cm,minimum height=1.0cm,align=center,font=\small},
  ahsrv/.style={draw=ArrowheadPurple,thick,fill=ArrowheadPurple!8,rounded corners=6pt,minimum height=0.55cm,inner xsep=4pt,align=center,font=\scriptsize\ttfamily},
  ahcsm/.style={draw=ArrowheadTeal,thick,fill=ArrowheadTeal!8,rounded corners=6pt,minimum height=0.55cm,inner xsep=4pt,align=center,font=\scriptsize\ttfamily},
  ahdb/.style={draw=ArrowheadOrange,thick,fill=ArrowheadOrange!8,cylinder,shape border rotate=90,aspect=0.25,minimum height=0.9cm,align=center,font=\scriptsize},
  ahflow/.style={-{Stealth[length=2.2mm]},thick,ArrowheadGrey},
  ahlabel/.style={font=\scriptsize,fill=white,inner sep=1pt,align=center},
}

%% Special references
\newcommand{\fref}[1]{{\textcolor{ArrowheadBlue}{\hyperref[sec:functions:#1]{#1}}}}
\newcommand{\mref}[1]{{\textcolor{ArrowheadPurple}{\hyperref[sec:model:#1]{#1}}}}
\newcommand{\sref}[1]{{\textcolor{ArrowheadPurple}{#1}}}
\newrobustcmd\msubsection[2]{
  \addtocounter{subsection}{1}
  \addcontentsline{toc}{subsection}{\protect\numberline{\thesubsection}#1 \textcolor{ArrowheadPurple}{#2}}
  \subsection*{\thesubsection\quad#1 \textcolor{ArrowheadPurple}{#2}}
  \label{sec:model:#2}
}

\begin{document}

%% Arrowhead Document Properties
\ArrowheadTitle{WorkflowConfigurator}

\ArrowheadType{System Description}
\ArrowheadTypeShort{SysD}
\ArrowheadVersion{5.4.0}
\ArrowheadDate{2026-11-24}
\ArrowheadAuthor{Documentation maintenance}
\ArrowheadStatus{REVIEW DRAFT}
\ArrowheadContact{Repository maintainers}
\ArrowheadFooter{\href{www.arrowhead.eu}{www.arrowhead.eu}}
\ArrowheadSetup
%%

%% Front Page
\begin{center}
  \vspace*{1cm}
  \huge{\arrowtitle}

  \vspace*{0.2cm}
  \LARGE{\arrowtype}
  \vspace*{1cm}

  
  \vspace*{\fill}

  

  \vspace*{1cm}
  \vspace*{\fill}

  % Front Page Abstract
  \begin{abstract}
    This document describes the WorkflowConfigurator, an Eclipse Arrowhead support system that lets an engineer draw workflows with conditions between the microservices of configured local clouds, and validates and stores them in the WorkflowManager. Documentation edition 5.4.0; system version 5.4.0 (framework), service versions 1.0.0.
  \end{abstract}

  \vspace*{1cm}
\end{center}

\newpage
%%

%% Table of Contents
\tableofcontents
\newpage
%%

\section{Overview}
\label{sec:overview}
The WorkflowConfigurator (WfC) is the design-time tool of the workflow systems. It offers a graphical
editor in a browser console in which an engineer selects a configured local cloud, sees the microservices and the
WorkflowExecutors' capabilities in that cloud, and draws a workflow: actions, steps, transitions with conditions,
dependencies, inputs and policies. The editor checks the drawing continuously, explains each finding with help on how
to fix it, and stores the workflow as a recipe in the WorkflowManager when it validates. The same operations are
offered as the Arrowhead service \texttt{workflow\allowbreak{}Configuration} for other engineering tools.

\begin{figure}[ht!]
  \centering
  \resizebox{\textwidth}{!}{
\begin{tikzpicture}[node distance=4mm]
\node[ahsys,minimum height=6.00cm,minimum width=3.6cm] (sys) at (0,0) {WorkflowConfigurator\\[-1pt]{\normalfont\scriptsize $\ll$system$\gg$}};
\node[ahsrv,anchor=west] (p0) at (3.4,2.62) {workflowConfiguration};
\draw[thick,ArrowheadPurple] (1.8,2.62) -- (p0.west);
\fill[ArrowheadPurple] (1.8,2.62) circle (1.4pt);
\node[ahsrv,anchor=west] (p1) at (3.4,1.88) {generalManagement};
\draw[thick,ArrowheadPurple] (1.8,1.88) -- (p1.west);
\fill[ArrowheadPurple] (1.8,1.88) circle (1.4pt);
\node[ahcsm,anchor=east] (c0) at (-3.4,2.62) {serviceDiscovery\\[-2pt]{\tiny\rmfamily ServiceRegistry}};
\draw[thick,ArrowheadTeal,-{Stealth[length=1.6mm]}] (c0.east) -- (-1.8,2.62);
\node[ahcsm,anchor=east] (c1) at (-3.4,1.88) {serviceRegistryManagement\\[-2pt]{\tiny\rmfamily ServiceRegistry}};
\draw[thick,ArrowheadTeal,-{Stealth[length=1.6mm]}] (c1.east) -- (-1.8,1.88);
\node[ahcsm,anchor=east] (c2) at (-3.4,1.12) {identity\\[-2pt]{\tiny\rmfamily Authentication}};
\draw[thick,ArrowheadTeal,-{Stealth[length=1.6mm]}] (c2.east) -- (-1.8,1.12);
\node[ahcsm,anchor=east] (c3) at (-3.4,0.38) {authorization\\[-2pt]{\tiny\rmfamily ConsumerAuthorization}};
\draw[thick,ArrowheadTeal,-{Stealth[length=1.6mm]}] (c3.east) -- (-1.8,0.38);
\node[ahcsm,anchor=east] (c4) at (-3.4,-0.38) {workflowRecipe\\[-2pt]{\tiny\rmfamily WorkflowManager}};
\draw[thick,ArrowheadTeal,-{Stealth[length=1.6mm]}] (c4.east) -- (-1.8,-0.38);
\node[ahcsm,anchor=east] (c5) at (-3.4,-1.12) {localCloudConfiguration\\[-2pt]{\tiny\rmfamily LocalCloudConfigurator}};
\draw[thick,ArrowheadTeal,-{Stealth[length=1.6mm]}] (c5.east) -- (-1.8,-1.12);
\node[ahcsm,anchor=east] (c6) at (-3.4,-1.88) {soLocalCloudConfiguration\\[-2pt]{\tiny\rmfamily SoLocalCloudConfigurator}};
\draw[thick,ArrowheadTeal,-{Stealth[length=1.6mm]}] (c6.east) -- (-1.8,-1.88);
\node[ahcsm,anchor=east] (c7) at (-3.4,-2.62) {sysMonitor (AMQP)\\[-2pt]{\tiny\rmfamily SysMonitor broker}};
\draw[thick,ArrowheadTeal,-{Stealth[length=1.6mm]}] (c7.east) -- (-1.8,-2.62);
\node[font=\scriptsize\bfseries,ArrowheadPurple] at (5.0,3.52) {produced};
\node[font=\scriptsize\bfseries,ArrowheadTeal] at (-5.0,3.52) {consumed};
\end{tikzpicture}
  }
  \caption{SysML block view of the WorkflowConfigurator: produced services (right) and consumed services with their providers (left).}
  \label{fig:block}
\end{figure}

The rest of this document is organized as follows.
In Section~\ref{sec:prior_art}, we reference major prior art capabilities of the system.
In Section~\ref{sec:use}, we describe the intended usage of the system.
In Section~\ref{sec:properties}, we describe fundamental properties provided by the system.
In Section~\ref{sec:delimitations}, we describe delimitations of capabilities of the system.
In Section~\ref{sec:services}, we describe the abstract service functions consumed or produced by the system.
In Section~\ref{sec:security}, we describe the security capabilities of the system.

\newpage
\subsection{Significant Prior Art}
\label{sec:prior_art}
Garcia Represa's workflow selection and implementation phases~\cite{represa2022}; the intent and
validator pipeline of the Arrowhead workflow kits; the LocalCloudConfigurator and SoLocalCloudConfigurator of this
workspace, whose designs the configurator lists.

\subsection{How This System Is Meant to Be Used}
\label{sec:use}
An operator opens the console (loopback, password login), picks the running cloud, a LocalCloudConfigurator
design or a member cloud of a System of LocalClouds design, and draws the workflow in the state-machine or sequence
view. Conditions are built with a condition builder (or typed in the guard grammar), inputs from templates referencing
parameters and results. \emph{Validate} asks the WorkflowManager; \emph{Save} registers the recipe; \emph{Open} loads
stored recipes. Import and export use the recipe JSON.

\subsection{System functionalities and properties}
\label{sec:properties}
\subsubsection{Functional properties of the system}
\begin{itemize}
\item Cloud directory: the live cloud from the ServiceRegistry (\texttt{service\allowbreak{}Registry\allowbreak{}Management} system-query and
service-query, with a discovery fallback), LocalCloudConfigurator designs and SoLC member clouds; each source optional.
\item Service and executor listing with operations, interfaces and capabilities; \texttt{sysMonitor} and
\texttt{general\allowbreak{}Management} hidden.
\item Graphical editor (static web application): state-machine and sequence views, inspector, condition builder,
undo, arrange with label placement, import/export, local checks with the validator's codes and per-finding help.
\item Forwarding of validate, store, list and get to the WorkflowManager's \texttt{workflow\allowbreak{}Recipe}; the list asks for
summaries, and a recipe document is read only when it is opened.
\item Retry policy in the step inspector: attempts, reaction to ERROR (\texttt{rebind} or \texttt{retry}) and backoff.
\end{itemize}

\subsubsection{Configuration of system properties}
The system is configured with Spring properties; the main ones are listed in the SysDD. The workflow-specific behaviour
is configured under the prefix \texttt{workflow.}; the security profile (TLS, authentication, authorization) is the
common Arrowhead~5 profile.

\subsubsection{Data stored by the system}
Nothing durable: recipes are stored only in the WorkflowManager. The log table is kept in an in-memory H2 database;
console sessions are kept in memory.

\subsubsection{Non functional properties}
\begin{itemize}
\item \emph{Security:} see Section~\ref{sec:security}; the console is a separate connector with its own authentication.
\item \emph{Latency:} each listing makes one to three calls to the ServiceRegistry or the configurators.
\item \emph{Energy and power saving:} idle between requests.
\end{itemize}

\subsubsection{Stateful or stateless}
Stateless apart from console login sessions.

\subsection{Important Delimitations}
\label{sec:delimitations}
\begin{itemize}
\item Services of a cloud other than the live one cannot be listed (its ServiceRegistry is not reachable); the editor lets
the engineer name service definitions.
\item The console is meant for loopback access; exposing it requires an explicit \texttt{console.\allowbreak{}address} and allowed
hosts.
\item The configurator does not start sessions; that is \texttt{workflow\allowbreak{}Instance} of the WorkflowManager.
\end{itemize}

\newpage
\section{Services}
\label{sec:services}
This section lists the produced and consumed services. Each produced service is specified by an SD and an HTTP/TLS/JSON
IDD in this documentation set.

\subsection{Produced services}

{\small
\begin{longtable}{|p{3.8cm}|p{6.2cm}|p{5.8cm}|}

\hline
\rowcolor{gray!33} Service & Operations & Documents \\ \hline
\endhead
\texttt{workflow\allowbreak{}Configuration} & list-clouds, list-services, list-executors, list-recipes, get-recipe, validate-recipe, store-recipe & \texttt{Workflow\allowbreak{}Configurator-\allowbreak{}v5.\allowbreak{}4.\allowbreak{}0-\allowbreak{}workflow\allowbreak{}Configuration-\allowbreak{}SD}, \texttt{-HTTP-IDD} \\ \hline
\texttt{general\allowbreak{}Management} & get-log, get-config & \texttt{Workflow-\allowbreak{}v5.\allowbreak{}4.\allowbreak{}0-\allowbreak{}general\allowbreak{}Management-\allowbreak{}SD}, \texttt{-HTTP-IDD} \\ \hline
\end{longtable}
}


\subsection{Consumed services}
All three workflow systems use the core systems the same way: at start-up they log in at the Authentication
system (\texttt{identity}) and register their services in the ServiceRegistry (\texttt{service\allowbreak{}Discovery.\allowbreak{}register});
on every incoming request they verify the caller's identity token (\texttt{identity.\allowbreak{}identity-\allowbreak{}verify}) and, with
\texttt{service.\allowbreak{}authorization.\allowbreak{}enabled}, the caller's grant (\texttt{authorization.\allowbreak{}verify} of ConsumerAuthorization);
at shutdown they revoke their services and log out.

{\small
\begin{longtable}{|p{3.4cm}|p{3.2cm}|p{3.6cm}|p{5.6cm}|}

\hline
\rowcolor{gray!33} Service & Provider & Operations & Purpose \\ \hline
\endhead
\texttt{service\allowbreak{}Discovery} & ServiceRegistry & register, revoke, lookup & Registration; executor and fallback discovery. \\ \hline
\texttt{service\allowbreak{}Registry\allowbreak{}Management} & ServiceRegistry & system-query, service-query & Systems and services of the live cloud (needs a grant and the \texttt{authorization} management policy). \\ \hline
\texttt{identity} & Authentication & login, logout, identity-verify & Own identity; verification of callers. \\ \hline
\texttt{authorization} & ConsumerAuthorization & verify & Provider-side authorization. \\ \hline
\texttt{workflow\allowbreak{}Recipe} & WorkflowManager & validate, register, query, get & Validation and storage. \\ \hline
\texttt{local\allowbreak{}Cloud\allowbreak{}Configuration} & LocalCloudConfigurator & catalog, list-configurations & Local-cloud designs. \\ \hline
\texttt{so\allowbreak{}Local\allowbreak{}Cloud\allowbreak{}Configuration} & So\allowbreak{}Local\allowbreak{}Cloud\allowbreak{}Configurator & list-designs, get-design & SoLC member clouds (when \texttt{workflow.\allowbreak{}configurator.\allowbreak{}solc.\allowbreak{}enabled}). \\ \hline
\texttt{sysMonitor (AMQP)} & SysMonitor broker & publish & When \texttt{sysmonitor.\allowbreak{}enabled}. \\ \hline
\end{longtable}
}


\newpage
\section{Security}
\label{sec:security}
The WorkflowConfigurator runs in Arrowhead secure mode in every cloud deployment; an insecure mode exists only for development and
unit tests. It handles only Arrowhead-compliant X.509 certificates of its local cloud; others are refused during the
TLS handshake.

\subsection{Security model}
\begin{itemize}
\item \textbf{Secure mode only in deployment.} The deployment configuration (\texttt{deployment/\allowbreak{}*.\allowbreak{}properties}) enables
TLS with client authentication (\texttt{server.\allowbreak{}ssl.\allowbreak{}client-\allowbreak{}auth=\allowbreak{}need}, TLS 1.3 and 1.2), outsourced authentication with
identity tokens bound to the client certificate, interface policy \texttt{CERT\_AUTH} and the CRL check. The packaged
defaults without TLS exist for development and unit tests (\texttt{workflow.\allowbreak{}standalone}) and must not be used in a
cloud.
\item \textbf{Certificates.} The system has an Arrowhead system certificate
\texttt{<System>.\allowbreak{}<cloud>.\allowbreak{}<organisation>.\allowbreak{}arrowhead.\allowbreak{}eu} issued by the local cloud CA, kept in a PKCS\#12 key store; the
trust store holds the cloud CA. Non-Arrowhead certificates are refused at the TLS layer.
\item \textbf{Provider-side authorization.} Every operation of the registered services (other than management) is
checked against ConsumerAuthorization (\texttt{authorization.\allowbreak{}verify}) for the calling system, the service definition
and the operation (route tables \texttt{Service\allowbreak{}Operations} and \texttt{Configurator\allowbreak{}Operations}). A missing grant gives
403.
\item \textbf{Management.} \texttt{general\allowbreak{}Management} and \texttt{workflow\allowbreak{}Management} are protected by
\texttt{management.\allowbreak{}policy=\allowbreak{}sysop-\allowbreak{}only}.
\item \textbf{Data protection.} All service traffic is encrypted by TLS. Each system is responsible for its own data:
its database runs inside its own process, so there is no database connection, port or account that another system could
use, and its data is reachable only through its services. Secrets (store and identity passwords, the WorkflowManager's
database password, the console password hash) come from container secrets (\texttt{configtree:/\allowbreak{}run/\allowbreak{}secrets/}), not from
the image.
\end{itemize}


\subsection{Console}
The browser console is a second connector (default \texttt{127.\allowbreak{}0.\allowbreak{}0.\allowbreak{}1:8485}). It authenticates the operator with a
password whose scrypt hash is configured as a secret, keeps the session in an HttpOnly cookie \texttt{wfc\_session},
checks the \texttt{Host} and \texttt{Origin} headers and locks out after five failed logins. A console session is not
an Arrowhead identity; the console cannot be reached on the Arrowhead connector, and the Arrowhead service cannot be
reached without mTLS and a grant.

\newpage
\bibliographystyle{IEEEtran}
\bibliography{bibliography}

\newpage

\section{Revision History}
\subsection{Amendments}

\noindent\begin{tabularx}{\textwidth}{| p{1cm} | p{2.2cm} | p{1.6cm} | X | p{3.6cm} |} \hline
\rowcolor{gray!33} No. & Date & Version & Subject of Amendments & Author \\ \hline
1 & 2026-11-23 & \arrowversion & Initial edition: system description of the WorkflowConfigurator, written from the Java reference implementation (ah5-Workflow-java). & Documentation maintenance \\ \hline
2 & 2026-11-24 & \arrowversion & Data protection; recipe summaries and the retry policy in the editor. & Documentation maintenance \\ \hline
\end{tabularx}

\subsection{Quality Assurance}

\noindent\begin{tabularx}{\textwidth}{| p{1cm} | p{2.2cm} | p{1.6cm} | X |} \hline
\rowcolor{gray!33} No. & Date & Version & Approved by \\ \hline
1 & & \arrowversion & Pending review \\ \hline
\end{tabularx}

\end{document}
